% 2026-11-11 / v38
\section{混合cubicのpoint schemeとPoisson零点集合}
\label{sec:points-poisson-cubic}

非可換のpoint schemeは、一般にはPoisson構造の零点集合と同じではない。
ここでは固定パラメータでの集合論的な記述を与え、h=0の極限、
中心三次商、形式復元先のcubicを区別する。

\subsection{混合中心三次商}
\[
F_1=(x^3+y^3+z^3)/3+u^3+v^3+uxy,
\]
\[
[y,z]=h(x^2+uy),\quad[z,x]=h(y^2+ux),\quad[x,y]=hz^2,
\quad[u,-]=[v,-]=0.
\]
この関係で定まる周囲の代数を$A_h$とし、
第\ref{mt:sec:family}節で構成した中心元$\Omega_h$による商を
$B_h=A_h/(\Omega_h)$とする。標数零の代数閉体において
$h\ne0$、$h^6\ne-27$とする。点モジュールはHilbert級数
$(1-t)^{-1}$の次数付き巡回右加群とし、標準の射影列$(p_i)_{i\ge0}$で表す。
ここでは集合論的なpoint parametrizationを扱い、
truncated point schemeの普遍族や非被約構造は同定しない。

\begin{proposition}[混合族のpoint moduleの台]\label{prop:psp-point-support}
上のパラメータでは、$B_h$の点モジュールは、その初期射影点によって、
次の二種類に分類される。
\[
\begin{split}
E_h&=\{u=v=0,\ x^3+y^3+z^3+h^2xyz=0\}
             \subset\mathbb P^4,\\
\Sigma&=\{z=0,\ x^2+uy=y^2+ux=0,\ F_1=0,\ u\ne0\}.
\end{split}
\]
$E_h$は滑らかなSklyanin楕円曲線で、$\Sigma$は12点からなる。
従ってpoint moduleの初期点の集合は
\[
|\operatorname{Point}(B_h)|=E_h\sqcup\Sigma.
\tag{PSP1}\label{eq:psp-support}
\]
$u$または$v$が非零の点モジュールはすべて定数列であり、
$E_h$上のものはSklyaninのtranslationで推移する。
\end{proposition}
\begin{proof}
各中心一次元$c=u,v$との交換関係を点モジュールの隣接座標
$p_i=(x_i,y_i,z_i,u_i,v_i)$、$p_{i+1}$へ多重線形化すると、
$c_ip_{i+1,a}=p_{i,a}c_{i+1}$（全一次生成元$a$）となる。
$c_i\ne0$なら$p_{i+1}=p_i$であり、
$c_{i+1}\ne0$も従う。前後の射影点にも同じ議論を適用するので、
中心座標がいずれかで非零の場合は全て定数列である。

定数列では交換子が零となるため、$h\ne0$から
\[
z^2=x^2+uy=y^2+ux=0
\]
を得る。点への中心元の作用は非可換語を同じスカラーに置く通常の代入である。
$z=0$を用いるとその値は$F_1$に一致する。
$u=0$なら$x=y=z=0$、$F_1=v^3\ne0$で不可能。
$u=1$に規格化すると、
$x=y=0$の場合は$v^3=-1$で三点。
$x,y\ne0$の場合は$xy=1$、$x^3=-1$、$y=-x^2$となり、
$v^3=-4/3$なので九点。いずれも定数列を定義する
（全関係式と中心元のスカラー代入が零）。

$u=v=0$の点列は三生成元のSklyanin代数
\[
S_h=\mathbb C\langle x,y,z\rangle/
(yz-zy-hx^2,zx-xz-hy^2,xy-yx-hz^2)
\]
の点列と同じである。この代数はSklyaninパラメータ
$(a,b,c)=(1,-1,-h)$を持つ。$h^6\ne-27$では
Sklyaninのpoint schemeが滑らかな楕円曲線
\[
abc(x^3+y^3+z^3)-(a^3+b^3+c^3)xyz=0
\]
であり、代入すると$E_h$を得る
\cite{PSPSmithTate,PSPDeLaet}.
この代数の中心三次元は全点モジュールを消す。
本研究の$\Omega_h\bmod(u,v)$はその非零スカラー倍なので、
すべての$S_h$点モジュールが$B_h$へ降りる。
中心座標が常に零の場合に他の点列はない。以上を合わせる。
\end{proof}

\begin{proposition}[Poisson零点との比較]\label{prop:psp-comparison}
$h=0$における$cubic$ $X_0=V(F_1)$のPoissonテンソル
\[
\{y,z\}=x^2+uy,\quad\{z,x\}=y^2+ux,\quad\{x,y\}=z^2
\]
の零点集合は台として
\[
|D_0(\pi_{X_0})|=E_0\sqcup\Sigma,
\quad E_0=\{u=v=0,\ x^3+y^3+z^3=0\}.
\tag{PSP2}\label{eq:psp-poisson}
\]
ここで臨界点部分の零点スキームは$24$の長さをもち、
$\Sigma$の各点で横断長2である。
一般の$h\ne0$では$E_h\ne E_0$であり、
$\operatorname{Point}(B_h)$は$X_0$に含まれさえしない。
一方、既存の有限対角から得た形式復元候補$G_h+v^3$について、
その$u=v=0$制限の零軌跡は$E_h$に等しい。
\end{proposition}
\begin{proof}
$E_0$で各活性JacobianのEulerベクトルが像に属するため、
射影化された双ベクトルは零になる。
$u=v=0$以外では中心方向が非零なので、射影化で階数は減らず、
Poisson零点の方程式は
\[
x^2+uy=y^2+ux=z^2=F_1=0.
\]
これは$\Sigma$の台を定める。
$u=1$上で$x^2+y=y^2+x=0$の解は
$(x,y)=(0,0)$または$(x^3,y)=(-1,-x^2)$で四個であり、
$x,y$方向のJacobian行列式はそれぞれ$-1$、$3$で非零。
各点で$v^3=-1$または$-4/3$なので$v$方向も単純である。
従って$z^2$が残る唯一の厚化で、12点それぞれ長さ2、合計24である。

$E_h$と$E_0$は$h\ne0$なら異なる三次曲線である。
例えば$h=1$では$p=[1:1:-1:0:0]$が$E_1$に属すが
$F_1(p)=1/3\ne0$。一般の$h$でも$xyz\ne0$の
$E_h$上の点は$F_1=-h^2xyz/3\ne0$。
最後に既存候補の$u=v=0$制限は
$(x^3+y^3+z^3+h^2xyz)/3$であるから一致する。
\end{proof}

\begin{remark}[埋め込みとスキーム論的分類]
$h=0$で$B_0$は可換cubic $X_0$の斉次座標環なので、
そのpoint schemeは$X_0$全体であり$D_0(\pi)$と等しくない。
$E_h\to E_0$という曲線の極限からpoint scheme全体の
flatな特殊化を結論してはならない。
本節の集合論的分類は次節でスキームとして精密化する。
孤立12点は各々長さ2の非被約点でPoisson零点の孤立厚化に
抽象的には同型だが、射影空間内の接方向は一般に異なる。
三次元Sklyaninのshiftは楕円曲線上の非自明なtranslationであり、
Poisson零点上の定数列とは異なる非可換動力学を保つ。
\end{remark}
